\scanpage{36}
\mspage{14}
$\pi_g$ is the smallest cardinal s.th. source and target of $g$ are $\pi$-accessible.

\uline{Example} If the objects \unsure{occuring|occurring} as sources and targets \deleted{of} $g$'s in $\Phi_0$ are ``of finite presentation'', i.e. $\mathrm{Hom}(A,-)$ commutes with any filtering direct limit, we get $\pi=0$. In this case we stop the induction with $\alpha_0=$ first infinite ordinal \inserted{$\omega$} \deleted{the first}. In case the objects $A,B$ are not of finite presentation (\inserted{\unsure{we assume|we can take|we mean}} still they are accessible --- we take the smallest cardinal such that $\mathrm{Hom}(A,?)$ commutes with direct limits indexed by \deleted{ordi} ordered sets \deleted{great} \inserted{large} w.r.t.\ $\pi$. Let \inserted{then} $\alpha_0$ be the first ordinal s.th. the \unsure{interval|ordset} $I_{\alpha_0}=\{\text{all ordinals }\alpha\text{ with }\alpha<\alpha_0\}$ be \deleted{great} large wrt $\pi$. We take the factorization

\diagram{\begin{tikzcd}[ampersand replacement=\&,row sep=small,column sep=large]
X \arrow[r,"i_{\alpha_0}(f)=i(f)"] \arrow[d] \& X_{\alpha_0}(f)=\deleted{\unsure{Z|\Sigma}}X(f) \arrow[dl,"\Sigma_{\alpha_0}(f)=p(f)"]\\
Y\&
\end{tikzcd}}
and we get thus $i(f)\in\Phi$, $p(f)\in\Psi$, provided we make the following assumptions;

d) The \deleted{\unsure{elem|objects}} source and target objects of the arrows of $\Phi_0$ are \uline{accessible} (ok if $M$ a topos etc ---) and, if $\pi_0$ is the cardinality of the smallest ordinal number \inserted{$\alpha_0$} s.th. for any such $A$, $\mathrm{Hom}(A,-)$ commutes to the \deleted{\unsure{f}} direct limit of type $I_{\alpha_0}$, \deleted{\unsure{we|we assume}} and \deleted{$\pi=\sup(\pi$} $\pi=\sup(\pi_0,\overset{\inserted{\text{rather, }\delta}}{\overbrace{\mathrm{card}\,\Phi_0}})$, we assume that for any filtering direct system with indexing set $I$ of cardinality $\leq\pi$, and transition maps in $\Phi$, the direct limit exists and $\deleted{\unsure{X_i}} Z_i\to Z$ are equally in $\Phi$.

Thus we get, not only factorization $f=pi$ with $i\in\Phi$, $p\in\Psi$, but even factorization $\deleted{\unsure{\Sigma|Z}}X(f)$ \uline{functorial} with regard to variable $f$.
% archive sheet number: 35
